% Part: applied-modal-logics
% Chapter: temporal-logic
% Section: temporal-logic-semantics

\documentclass[../../../include/open-logic-section]{subfiles}

\begin{document}

\olfileid{aml}{tl}{sem}

\olsection{دلالات المنطق الزمني}

\begin{defn}
تحوي اللغة الأساسية للمنطق الزمني ما يأتي:
\begin{enumerate}
  \tagitem{prvFalse}{الثابت القضي الدال على !!{falsity}~$\lfalse$.}{}
  \tagitem{prvTrue}{الثابت القضي الدال على !!{truth}~$\ltrue$.}{}
  \item مجموعة !!{denumerable} من ال!!{propositional variable}s:
    $\Obj p_0$، $\Obj p_1$، $\Obj p_2$، \dots
  \item الروابط القضية: \startycommalist
  \iftag{prvNot}{\ycomma $\lnot$ (النفي)}{}%
  \iftag{prvAnd}{\ycomma $\land$ (العطف)}{}%
  \iftag{prvOr}{\ycomma $\lor$ (الفصل)}{}%
  \iftag{prvIf}{\ycomma $\lif$ (!!{conditional})}{}%
  \iftag{prvIff}{\ycomma $\liff$ (!!{biconditional})}{}.
  \item مؤثرا الماضي~$\Ptemp$ و$\Htemp$.
  \item مؤثرا المستقبل~$\Ftemp$ و$\Gtemp$.
\end{enumerate}
\end{defn}

سنناقش لاحقًا إمكان إضافة أنواع أخرى من المؤثرات الجهوية.

\begin{defn}
تُعرّف \emph{!!^{formula}s} اللغة الزمنية استقرائيًا كما يأتي:
\begin{enumerate}
\tagitem{prvFalse}{$\lfalse$ !!{formula} ذرية.}{}

\tagitem{prvTrue}{$\ltrue$ !!{formula} ذرية.}{}

\item كل متغير قضي~$\Obj p_i$ هو !!{formula} (ذرية).

\tagitem{prvNot}{إذا كانت~$!A$ !!a{formula}، كانت~$\lnot !A$
  !!a{formula}.}{}

\tagitem{prvAnd}{إذا كانت~$!A$ و$!B$ !!{formula}s، كانت
  $(!A \land !B)$ !!a{formula}.}{}

\tagitem{prvOr}{إذا كانت~$!A$ و$!B$ !!{formula}s، كانت
  $(!A \lor !B)$ !!a{formula}.}{}

\tagitem{prvIf}{إذا كانت~$!A$ و$!B$ !!{formula}s، كانت
  $(!A \lif !B)$ !!a{formula}.}{}

\tagitem{prvIff}{إذا كانت~$!A$ و$!B$ !!{formula}s، كانت
  $(!A \liff !B)$ !!a{formula}.}{}

\item إذا كانت~$!A$ !!a{formula}، فإن~$\Ptemp !A$ و$\Htemp !A$
  و$\Ftemp !A$ و$\Gtemp !A$ كلها !!{formula}s.

\tagitem{limitClause}{لا شيء غير ذلك !!a{formula}.}{}
\end{enumerate}
\end{defn}

تُعطى دلالات المنطق الزمني بواسطة نماذج علائقية، كما في سائر أنواع المنطق
المقصدي.

\begin{defn}
  \emph{نموذج} اللغة الزمنية ثلاثي
  $\mModel{M}=\tuple{T,\prec,V}$، حيث:
  \begin{enumerate}
  \item $T$ مجموعة غير فارغة، تُفسر عناصرها بوصفها نقاطًا في الزمن.
  \item $\prec$ علاقة ثنائية على~$T$.
  \item $V$ دالة تسند إلى كل~!!{propositional
    variable}~$p$ مجموعة~$V(p)$ من النقاط الزمنية.
  \end{enumerate}
  وحين يكون~$t \prec t'$، نقول إن~$t$ \emph{يسبق}~$t'$.
  وحين يكون~$t \in V(p)$، نقول إن~$p$ \emph{صادق عند}~$t$.
\end{defn}

ستلاحظ أننا لا نفرض الآن أي شروط على علاقة السبق~$\prec$. وهذا يعني أنه
لا توجد حاليًا قيود على بنية نماذجنا الزمنية؛ فيمكن أن تكون لدينا نماذج
يكون الزمن فيها خطيًا أو متفرعًا أو دائريًا، أو ذا أي بنية كانت.

وكما في المنطق الجهوي النظامي، يحدد كل نموذج زمني أي ال!!{formula}s تُعد
صادقة عند أي نقاط فيه. ونستعمل الترميز نفسه، «النموذج~$\mModel{M}$ يصدق
!!{formula}~$!A$ عند النقطة~$t$»، للمفهوم الأساسي في الدلالات العلائقية.
وتُعرّف العلاقة استقرائيًا، وهي مطابقة لحالة المنطق الجهوي النظامي في جميع
المؤثرات غير الجهوية.

\begin{defn}\ollabel{defn:tmodels}
  يُعرّف \emph{صدق !!a{formula}~$!A$ عند~$t$} في~$\mModel M$، ورمزه
  $\mSat{M}{!A}[t]$، استقرائيًا كما يأتي:
  \begin{enumerate}
  \tagitem{prvFalse}{\indcase{!A}{\lfalse}{لا يتحقق أبدًا
    $\mSat{M}{\lfalse}[t]$}.}{}
  \tagitem{prvTrue}{\indcase{!A}{\ltrue}{يتحقق دائمًا
    $\mSat{M}{\ltrue}[t]$}.}{}
  \item $\mSat{M}{p}[t]$ إذا وفقط إذا كان~$t \in V(p)$.
  \tagitem{prvNot}{\indcase{!A}{\lnot !B}{$\mSat{M}{\indfrm}[t]$ إذا
    وفقط إذا كان~$\mSat/{M}{!B}[t]$}.}{}
  \tagitem{prvAnd}{\indcase{!A}{(!B \land !C)}{$\mSat{M}{\indfrm}[t]$
    إذا وفقط إذا كان~$\mSat{M}{!B}[t]$ و$\mSat{M}{!C}[t]$}.}{}
  \tagitem{prvOr}{\indcase{!A}{(!B \lor !C)}{$\mSat{M}{\indfrm}[t]$
    إذا وفقط إذا كان~$\mSat{M}{!B}[t]$ أو~$\mSat{M}{!C}[t]$}
    (أو كلاهما).}{}
  \tagitem{prvIf}{\indcase{!A}{(!B \lif !C)}{$\mSat{M}{\indfrm}[t]$
    إذا وفقط إذا كان~$\mSat/{M}{!B}[t]$ أو~$\mSat{M}{!C}[t]$}.}{}
  \tagitem{prvIff}{\indcase{!A}{(!B \liff !C)}{$\mSat{M}{\indfrm}[t]$
    إذا وفقط إذا تحقق كل من~$\mSat{M}{!B}[t]$ و$\mSat{M}{!C}[t]$، أو
    لم يتحقق أي من~$\mSat{M}{!B}[t]$ و$\mSat{M}{!C}[t]$}.}{}
  \item\ollabel{defn:sub:mmodels-p}
    \indcase{!A}{\Ptemp !B}{$\mSat{M}{\indfrm}[t]$ إذا وفقط إذا كان
    $\mSat{M}{!B}[t']$ لبعض~$t' \in T$ بحيث~$t' \prec t$}
  \item\ollabel{defn:sub:mmodels-h}
    \indcase{!A}{\Htemp !B}{$\mSat{M}{\indfrm}[t]$ إذا وفقط إذا كان
    $\mSat{M}{!B}[t']$ لكل~$t' \in T$ بحيث~$t' \prec t$}
  \item\ollabel{defn:sub:mmodels-f}
    \indcase{!A}{\Ftemp !B}{$\mSat{M}{\indfrm}[t]$ إذا وفقط إذا كان
    $\mSat{M}{!B}[t']$ لبعض~$t' \in T$ بحيث~$t \prec t'$}
  \item\ollabel{defn:sub:mmodels-g}
    \indcase{!A}{\Gtemp !B}{$\mSat{M}{\indfrm}[t]$ إذا وفقط إذا كان
    $\mSat{M}{!B}[t']$ لكل~$t' \in T$ بحيث~$t \prec t'$}
  \end{enumerate}
\end{defn}

قد يتضح لك من الدلالات أن المؤثرين~$\Ptemp$ و$\Htemp$ متقابلان، وكذلك
المؤثران~$\Ftemp$ و$\Gtemp$؛ بحيث يمكننا تعريف~$\Htemp !A$ بأنه
$\lnot \Ptemp \lnot !A$، وفعل الأمر نفسه بـ~$\Gtemp$ و$\Ftemp$.

\end{document}
